\documentclass[12pt]{exam}
\usepackage{amsmath,amssymb}
\usepackage[margin=0.5in]{geometry}
\usepackage{tikz}
\usepackage{listings}
\usepackage{inconsolata}
\usepackage[T1]{fontenc}
\lstset{basicstyle=\ttfamily}
\usepackage{soul}

\newcommand{\blank}[1]{\underline{\hspace*{#1}}}
\newcommand{\ds}{\displaystyle}
\newcommand{\on}{\operatorname}


\addpoints

\lstset{
    numbers=left,
    frame=l,
    basicstyle=\ttfamily\small,
    language=Python
}


\begin{document}
\pagestyle{empty}
\graphicspath{{/home/brian/Dropbox/HSC/Spring16/Math111/}}
\lstset{language=Python}
\lstset{columns=fixed}
\lstset{showstringspaces=false}

\subsection*{Midterm 1 Practice Exam   \hfill CS 261}

\begin{questions}

\question[8] Complete lines 5 and 7 in the following function. 
\begin{lstlisting}
def weeks_and_days(days):
    """Converts any number of days into weeks plus extra days.  
    For example, if the input is 30, it should print '4 week(s) and 2 day(s)'.
    """
    weeks = 

    days = 

    print(weeks, "week(s) and", days, "day(s).")
\end{lstlisting}

\question[12] For each of the following Python expressions, write down the output when the expression is evaluated using a Python interpreter. Write \textbf{error} if you think the expression will raise an error.

\begin{parts}
\part \lstinline{19 % 5}
\vfill

\part \lstinline{5 + 2 * 10}
\vfill

\part \lstinline{"HSC" * 3}
\vfill

\part \lstinline{str(1) + int(2)}
\vfill

\part \lstinline{13 // 4}
\vfill

\part \lstinline{1 + 2 = 3}
\vfill
\end{parts}

\question[8] What will the following function return when you call it? Be sure to give both the return value and its type.
\begin{lstlisting}
def mystery():
    x = 5.0
    y = int(x)
    z = y ** 2
    return z - x
\end{lstlisting}
\vfill

\question[12] Write a loop with an accumulator variable that takes a string \lstinline{s} and creates a new string where every letter or character in \lstinline{s} is repeated twice.  For example, if \lstinline{s = "Hello!"}, your loop should generate the string \lstinline{"HHeelllloo!!"}.
\vspace*{2.5in}
\vfill

\newpage

\question[10] Identify the type (str, int, float, bool, or list) of each expression.
\begin{parts}
\part \lstinline{["hello", "how", "are", "you?"]}
\bigskip

\part \lstinline{6 / 2}
\bigskip

\part \lstinline{5 < 6}
\bigskip

\part \lstinline{-1 + 4 // 5}
\bigskip

\part \lstinline{"4"}
\bigskip
\end{parts}



\question[8] Consider the following function. 

\begin{lstlisting}
def f(x):
    while x > 3:
        x = x + 1
    return x
\end{lstlisting}

\begin{parts}
\part What will happen if you call \lstinline{f(5)}? 
\vfill


\part What will happen if you call \lstinline{f(2)}? 
\vfill

\end{parts}

\question[12] A list is \emph{increasing} if every element in the list is greater than the element before it. The following function is supposed to return \lstinline{True} if a list is increasing and \lstinline{False} otherwise, but it has a bug.

%    """Supposed to return True if values is an increasing list. 
%    But there is a bug in the code.
%    """
\begin{lstlisting}
def is_increasing(values):
    for i in range(1, len(values)):
        if values[i-1] < values[i]:
            return True
        else:
            return False
\end{lstlisting}

\begin{parts}
\part What will \lstinline{is_increasing([2, 3, 5])} return?
\vfill

\part What will \lstinline{is_inscreasing([4, 8, 1, 5])} return? 
\vfill

\part What is wrong with the way the function \lstinline{is_increasing} is coded?
\vfill
\end{parts}

\question[4] What will the following code output?

\begin{lstlisting}
x = [1, 2, 3, 4]
print(x[2])
\end{lstlisting}
\vfill

\newpage
\question[8] Consider the following recursive function which accepts an integer input. 

\begin{lstlisting}
def recurse(n):
    if n <= 1:
        return 0
    else:
        return 1 + recurse(n // 2) 
\end{lstlisting}

\begin{parts}
\part What does \lstinline{recurse(-1)} return?
\vfill

\part What does \lstinline{recurse(10)} return?
\vfill

\end{parts}
\question[8] Determine the values of the following boolean expressions. 
\begin{parts}
\part \lstinline{5 < 6 and 2 + 2 == 4}
\vfill 

\part \lstinline{print(4 + 2) == 6}
\vfill 


\part \lstinline{"Fish"  in "All the fish in the sea."}
\vfill 

\part \lstinline{not (5 > 6 or 1+1 != 3)}
\vfill 

\end{parts}


\question[10] Trace the following function by making a table to show how the values of the variables \lstinline{t}, \lstinline{x}, and \lstinline{n} change when you compute \lstinline{result([1, 5, 4])}.  Be sure to indicate what the function returns at the end.


\begin{lstlisting}
def result(numbers):
    t = 0
    for x in numbers:
        n = 2 * x
        t = t + n - 1
    return t
\end{lstlisting}

\vspace*{2.5in}
\vfill
\end{questions}


\end{document}



\color{gray}
\question The following code is supposed to print the perfect squares 1, 4, 9, 16, $\ldots$ up to 100.  Circle the mistake in the code and explain why it doesn't do what it is supposed to. 
\begin{lstlisting}
n = 1
while n <= 10:
    square = n ** 2
    n + 1
    print(square)
\end{lstlisting}
\begin{solution}
The mistake is in the line \lstinline{n + 1}.  It never changes the value of \lstinline{n}, so the loop goes on forever.
\end{solution}
 
%
%\question 100 minutes is 1 hour and 40 minutes. How many hours and minutes are in 291 minutes?  You don't need to calculate the answer.  Just complete the following variable assignments with expressions that would calculate the answer.
%\lstset{language=Python}
%\begin{lstlisting}
%hours = 
%minutes = 
%\end{lstlisting}
%\begin{solution}
%\begin{lstlisting}
%hours = 291 // 60
%minutes = 291 % 60
%\end{lstlisting}
%
%\end{solution}

%\question For each of the following Python expressions, write down the output when the expression is evaluated using a Python interpreter. Write \textbf{error} if you think the expression will raise an error.
%
%\begin{parts}
%\part \lstinline{15 % 3}
%\begin{solution}
%\lstinline{0}
%\end{solution}
%\bigskip
%
%\part \lstinline{5 + 3 * 4 == 100}
%\begin{solution}
%False
%\end{solution}
%\bigskip
%
%\part \lstinline{"HSC" + "2024"}
%\begin{solution}
%\lstinline{"HSC2024"}
%\end{solution}
%\bigskip
%
%\part \lstinline{100+"1"}
%\begin{solution}
%\textbf{error}
%\end{solution}
%\bigskip
%
%\part \lstinline{6 / 2}
%\begin{solution}
%\lstinline{3.0}
%\end{solution}
%\bigskip
%\end{parts}

\question Suppose that we type the following assignments and expressions in a Python
shell in the given order. 

\begin{verbatim}
    >>> a = 10
    >>> b = 20
    >>> c = a + b
\end{verbatim}

\begin{parts}
\part What will Python output if we enter the following (after those first three statements)?

\begin{verbatim}
>>> a + 5
\end{verbatim}
\begin{solution}
15
\end{solution}
\bigskip

\part What will Python output if we enter this statement (after the one from part (a))?
\begin{verbatim}
>>> a + b
\end{verbatim}
\begin{solution}
30
\end{solution}

\bigskip

\part What will Python output when we enter these two statements next?
\begin{verbatim}
>>> b = b + 1
>>> c
\end{verbatim}
\begin{solution}
30
\end{solution}
\bigskip

\part What about if we enter this last?
\begin{verbatim}
>>> a + b
\end{verbatim}
\begin{solution}
31
\end{solution}

\end{parts} 

\newpage
\question Consider the following function.  What will it return if you call \lstinline{mystery(5)}?

\begin{lstlisting}
def mystery(n):
    a = n 
    a + 5
    if a > 10:
        return a
    elif a == 10:
        return 100
    else:
        return n
\end{lstlisting}
\begin{solution}
5
\end{solution}

%\question Write a function called \lstinline{average3} that inputs any three numbers and returns their average.  
%\begin{solution}
%\begin{lstlisting}
%def average3(a,b,c):
%    return (a + b + c) / 3
%\end{lstlisting}
%\end{solution}
%\vfill

\question Consider the simple function below.

\begin{lstlisting}
def twice(n):
    print(2 * n)
\end{lstlisting}

Why do we get \lstinline{False} if we enter the following into the shell? Explain why we get the output below. 

\begin{verbatim}
    >>> twice(5) == 10
    10
    False
\end{verbatim}

\begin{solution}
The function prints 10, but \lstinline{twice(5)} is not equal to 10 because the function is returning nothing.  
\end{solution}
\vfill

\question Describe in words what the following function will do when you call it.

\begin{lstlisting}
def loop_function():
    for i in range(100):
        if i % 3 == 0:
            print(i)
\end{lstlisting}

\begin{solution}
This function will print the integers that are multiples of 3, starting at 0, until 99. 
\end{solution}

\question Identify the type of the value (str, int, float, bool, or list) for each of the following expressions.
\begin{parts}
\part \lstinline{"hello" * 5}
\begin{solution}
\lstinline{str}
\end{solution}
\bigskip

\part \lstinline{ 7 // 2}
\begin{solution}
\lstinline{int}
\end{solution}
\bigskip

\part \lstinline{"hello" == 5}
\begin{solution}
\lstinline{bool}
\end{solution}
\bigskip

\part \lstinline{("He" in "Hello") or ("Yes" in "No")}
\begin{solution}
\lstinline{bool}
\end{solution}
\bigskip

\part \lstinline{3 + 4 / 7}
\begin{solution}
\lstinline{float}
\end{solution}
\bigskip

\end{parts}

\question Determine the values of the variables $x$, $y$, and $z$ after the following lines of code are run. 


\begin{lstlisting}
x = 5
y = x
z = 2 * y / 2
x += 4
y = (y - 1) * 2
x = (x == z + 4)
\end{lstlisting}

\begin{solution}
$y$ is 8, $z$ is 5.0, and $x$ is \lstinline{True}. 
\end{solution}

\question Consider the following recursive function: 


\begin{lstlisting}
def recurse(n, s):
    if n == 0:
        return s
    else: 
        return recurse(n - 1, s + n)
\end{lstlisting}
\begin{parts}
\part What will \lstinline{recurse(3, 0)} return?
\begin{solution}
6
\end{solution}
\bigskip

\part What will happen if you call \lstinline{recurse(-1, 0)}?
\begin{solution}
It will enter an infinite recursion because \lstinline{n} is never 0. 
\end{solution}
\bigskip
\end{parts}



\end{questions}



\end{document} 


Certainly! Here are some example questions you could use for a midterm exam in an Introduction to Computer Science with Python class. The questions cover a range of topics, including basic syntax, data structures, algorithms, and problem-solving.

### Multiple Choice Questions

1. **Which of the following is the correct syntax to define a function in Python?**
   a) `function myFunction()`
   b) `def myFunction:`
   c) `def myFunction():`
   d) `myFunction() {}`

2. **What will the following code output?**
   ```python
   x = [1, 2, 3, 4]
   print(x[2])
   ```
   a) `1`  
   b) `2`  
   c) `3`  
   d) `4`

3. **What is the result of the following expression?**
   ```python
   3 * (4 + 2) - 5
   ```
   a) `13`  
   b) `18`  
   c) `10`  
   d) `16`

### Short Answer Questions

4. **Explain the difference between a list and a tuple in Python.**

5. **What is the purpose of the `if __name__ == "__main__":` statement in a Python program?**

### Coding Questions

6. **Write a Python function named `is_even` that takes an integer as an argument and returns `True` if the integer is even and `False` if it is odd.**

   ```python
   def is_even(n):
       # Your code here
   ```

7. **Given the following list: `numbers = [10, 20, 30, 40, 50]`, write a loop to print each number multiplied by 2.**

   ```python
   numbers = [10, 20, 30, 40, 50]
   # Your code here
   ```

### Problem-Solving Questions

8. **You have a string `"hello world"`. Write a Python function `count_vowels` that returns the number of vowels in the string.**

   ```python
   def count_vowels(s):
       # Your code here
   ```

9. **Write a Python program that generates a list of the first 10 Fibonacci numbers.**

   ```python
   def fibonacci(n):
       # Your code here
   ```

### True/False Questions

10. **True or False: In Python, all variables must be declared before they are used.**

11. **True or False: A Python list can contain elements of different data types.**

Feel free to modify these questions based on the specific content covered in your course!
